\section{The Decision Procedure}\label{sec:correct}

Section~\ref{sec:tableau} reduces satisfiability of $\varphi$ to whether $N_0$ lies in a successful
marking of $\TabGraph$. The procedure that decides this rests on reading the conditions of
\S\ref{ss:loop} as a removal operator $F$ on subsets of $\TabNode$. 
%itemjoin={{; }}, itemjoin*={{; and }}
For $A\subseteq\TabNode$, a node
$v\in A$ survives when
\begin{enumerate}[label={\footnotesize\textbf{(D\arabic*)}}, ref={\footnotesize\textbf{(D\arabic*)}},itemindent=2em,
                   ]
\item\label{D1} $v$ is not closed
\item\label{D2} $\ch(v)\cap A\neq\emptyset$ if $\kind(v)=\ORn$, and $\ch(v)\subseteq A$ if
$\kind(v)=\ANDn$
\item\label{D3} $v$ satisfies \ref{M1} and \ref{M2} with its paths taken inside $A$
\end{enumerate}
\noindent
that is, $F(A)=\{v\in A \mid \text{$v$ satisfies \ref{D1}--\ref{D3} on $A$} \}.$ 

$F$ is monotone by construction. Let $A\subseteq B\subseteq\TabNode$ and let $v\in F(A)$. Then
$v\in B$, and each of \ref{D1}--\ref{D3} still holds of $v$ on $B$: \ref{D1} does not mention the
set; a \ORn-node keeps in $B$ the child it had in $A$, and an \ANDn-node has all its children in
$A\subseteq B$; and a path inside $A$ is a path inside $B$, which settles \ref{D3} and \ref{D3}.
Hence $v\in F(B)$ and $F(A)\subseteq F(B)$.


%%%%%%%%% DEAD TEXT %%%%%%%%
% The same operator stated by what it deletes rather than what it keeps; kept in
% case the removal phrasing reads better in the (D3)/(D4) clauses.
%\S\ref{ss:loop} as an operator $F$ which removes from a set $A\subseteq\TabGraph$:
%\begin{enumerate*}[label={\textbf{(D\arabic*)}}, ref={\textbf{(D\arabic*)}},
%                   itemjoin={{; }}, itemjoin*={{; and }}]
%\item the closed nodes
%\item the \ORn-nodes with no child in $A$, and the \ANDn-nodes with a child outside $A$
%\item the nodes at which \ref{M1} fails when its paths are taken inside $A$
%\item the nodes at which \ref{M2} fails when its paths are taken inside $A$.
%\end{enumerate*}
%%%%%%%%% DEAD TEXT %%%%%%%%



%
%Let $A\subseteq B$ and $v\in F(A)$, so $v\in A$ and $v$ survives every clause on $A$. Then
%$v\in B$, and it survives them on $B$ as well: \ref{D1} does not mention the set; a \ORn-node
%keeps its child in $A\subseteq B$; an \ANDn-node has no child outside $A$, and a child outside
%$B$ would be one; and a path inside $A$ is a path inside $B$, which settles \ref{D3} and
%\ref{D4}. So $v\in F(B)$ and $F$ is monotone.

\begin{lemma}\label{prop:gfp}
\TabGraph has a greatest successful marking $\mathcal{M}^*$ and succeeds
iff $N_0\in\mathcal{M}^*$.
\end{lemma}
\begin{proof}
Since $F$ is monotone and $\TabGraph$ is finite by Lemma~\ref{lem:finiteness}, by Knaster--Tarski, $F$ has a greatest fixpoint $\nu F$. Every marking surviving \ref{D1}--\ref{D3} is successful by construction, hence $\nu F$ is the greatest successful marking. As a consequence, the tableau succeeds on $\varphi$ iff $N_0\in\nu F$.
\end{proof}


A procedure that naively iterates $F$ downwards from $\TabGraph$ reaches $\nu F$ in finitely many steps, but presupposes
the full construction of $\TabGraph$. Instead, Algorithm~\ref{alg:otf} decides $N_0\in\nu F$ on-the-fly, building only the part of $\TabGraph$ required. It works by expanding each node of \TabGraph until its children (line~\ref{ln:expand}). Lines~\ref{ln:seed}--\ref{ln:decide} take a strongly
connected component (SCC) $C$ of the expansion whose children outside $C$ are all either in $\Acc$
or $\Rej$, apply $F$ inside $C$, accept what survives and reject the rest (Line~\ref{ln:decide}). $\Acc$ is
passed to $F$ at line~\ref{ln:fix}, so a node of $C$ whose witness leaves $C$ through an accepted
child finds it continued there.
The loop exits once $N_0$ is decided, and $\varphi$ is
satisfiable exactly when $N_0$ ends in $\Acc$. Lines~\ref{ln:acc} and~\ref{ln:rej} are optimizations which allow pruning: a \ORn-node with an accepted child, or an \ANDn-node with
a rejected one, is decided on the spot, so only the components of $E\setminus (\Acc\cup\Rej)$ are ever computed.

\begin{algorithm}[t]
\footnotesize
\DontPrintSemicolon
\SetAlgoLined
\SetAlgoNoEnd
\KwIn{$\varphi$ in negation normal form}
\KwOut{\SAT{} or \UNSAT}
$E\gets\{N_0\}$; \quad $\Acc\gets\emptyset$; \quad $\Rej\gets\emptyset$\label{ln:init}\;
\Repeat{$N_0\in \Acc \cup \Rej$\label{ln:until}}{
  $\Acc\gets\Acc\cup\{v\in E\mid \kind(v)=\ORn,\ \ch(v)\cap\Acc\neq\emptyset\}$\label{ln:acc}\;
  $\Rej\gets\Rej\cup\{v\in E\mid \kind(v)=\ANDn,\ \ch(v)\cap\Rej\neq\emptyset\}$\label{ln:rej}\;
  \ForEach{component $C$ of $E\setminus (\Acc\cup\Rej)$ with $\ch(C)\setminus C\subseteq \Acc\cup\Rej$\label{ln:comp}}{
    $X\gets C$\label{ln:seed}\;
    \lRepeat{$X$ \textnormal{is unchanged}}{$X\gets F(X\cup\Acc)\cap C$\label{ln:fix}}
    $\Acc\gets\Acc\cup X$; \quad $\Rej\gets\Rej\cup(C\setminus X)$\label{ln:decide}\;
  }
  $E\gets E\cup\ch(E\setminus (\Acc\cup\Rej))$\tcp*{apply the rules of \S\ref{ss:rules}}\label{ln:expand}
}
\Return{\SAT{} \textnormal{if} $N_0\in\Acc$, \textnormal{else} \UNSAT}\label{ln:return}
\caption{Satisfiability check for TCTL formula $\varphi$}\label{alg:otf}
\end{algorithm}

Correctness of the procedure rests on three further facts: \ref{lemitem:pruning} pruning preserves $\nu F$; \ref{lemitem:de} a time-progressing cycle is found by inspecting resets and bounds, making \ref{D3} computable; and \ref{lemitem:fix} $F$ may be applied inside a single component, preserving the greatest fixpoint and allowing for an on-the-fly decision.

\begin{lemma}\label{lem:search}
$A\subseteq\mathcal M^{*}$ be a successful marking with $\mathcal M^{*}\cap(\ch(C)\setminus C)\subseteq A$.
\begin{enumerate}[label=\emph{(\roman*)}, ref={(\roman*)}]
\item\label{lemitem:pruning} If $v$ is a \ORn-node with a child $w\in\mathcal M^{*}$, or an
\ANDn-node with a child $w\notin\mathcal M^{*}$, then $v\in\mathcal M^{*}$ iff
$w\in\mathcal M^{*}$. Removing $v$'s other children leaves $\mathcal M^{*}$ unchanged.
\item\label{lemitem:de} $T$ contains a time-progressing cycle iff some node of $T$ is a
\ruleref{Delay}-successor and every clock is unbounded or reset in $T$.
\item\label{lemitem:fix} $\mathcal M^{*}\cap C$ is the largest $X\subseteq C$ with
$X\subseteq F(X\cup A)$.
\end{enumerate}
\end{lemma}
\begin{proof}
\ref{lemitem:pruning} For an \ANDn-node this is \ref{D2}. For a \ORn-node, $\mathcal M^{*}\cup\{v\}$
is a post-fixpoint of $F$: $v$ is not closed and spared by \ref{D2} because of $w$; and \ref{D3} hold at $v$ by prefixing the edge
$v\to w$ to the witnesses $w$ carries, a pursuit active at $v$ being either still active at $w$ or
discharged by that edge. Hence $v\in\nu F=\mathcal M^{*}$. The same prefixing gives the second
claim: a path using another child of $v$ arrives through $v$, so a node relying on it can take its
own path to $v$ and then $v$'s witness. \ref{lemitem:de} ($\Leftarrow$) Strong connectivity gives a cycle through the
\ruleref{Delay}-successor and a reset node of each clock $T$ resets, which is then
time-progressing. ($\Rightarrow$) By definition a time-progressing cycle elapses positive time, so it crosses a \ruleref{Delay}, since it is the only rule that elapses time, and every clock is either reset on it or extrapolated above $M$, hence unbounded. 
\ref{lemitem:fix} For $\subseteq$, $C$ being maximal a witness path of $\mathcal M^{*}$ from
$N\in C$ either stays in $C$ or leaves once through some $c\in A$, where it continues since $A$ is a
successful marking. For $\supseteq$, $X\cup A$ is then a successful marking too, its $A$-part by
monotonicity and its $X$-part by hypothesis, so $X\subseteq\mathcal M^{*}$ by
Proposition~\ref{prop:gfp}.
\end{proof}

\begin{theorem}\label{thm:alg}
Algorithm~\ref{alg:otf} terminates in $2^{\mathrm{poly}(n)}$ and returns \SAT{} iff $N_0\in\nu F$.
\end{theorem}



\begin{proof}
Let $D=\Acc\cup\Rej$ and $U=E\setminus D$. Each iteration enlarges $E$ or $D$, or the loop exits.
Suppose one enlarges $D$ by nothing; then $E$, $D$ and $U$ still carry the values they had on
entry when line~\ref{ln:expand} runs. If $U=\emptyset$ then $E=D$, so $N_0\in D$ and the loop
exits. Otherwise suppose line~\ref{ln:expand} adds nothing either, that is $\ch(U)\subseteq E$.
The condensation of the subgraph induced on $U$ is non-empty, finite and acyclic, so it has a sink
$C$: a maximal SCC of $U$ with $\ch(C)\cap U\subseteq C$. Take any $w\in\ch(C)\setminus C$. It
lies outside $U$ by sinkhood, and inside $E$ by that supposition, hence in $D$. So every exit of
$C$ is decided, $C$ meets the guard of line~\ref{ln:comp}, and line~\ref{ln:decide} adds the
non-empty $C$ to $D$, against the assumption that $D$ did not grow. Hence
$\ch(U)\not\subseteq E$, and line~\ref{ln:expand} enlarges $E$. Both sets are bounded by $|\TabNode|$, finite by Lemma~\ref{lem:finiteness}, so at most
$2|\TabNode|$ iterations enlarge either and the next one exits, necessarily with $N_0\in D$ as
that is the guard of line~\ref{ln:until}.



We show $\Acc\subseteq\mathcal M^{*}$ and $\Rej\cap\mathcal M^{*}=\emptyset$ by induction on the
order in which nodes are accepted or rejected, each such step resting on those already taken.
Lines~\ref{ln:acc} and~\ref{ln:rej} are Lemma~\ref{lem:search}\ref{lemitem:pruning}, which also
justifies leaving the remaining children unbuilt. Line~\ref{ln:comp} fires only once every child
of $C$ outside $C$ has been accepted or rejected, so by induction
$\Acc\cap\ch(C)=\mathcal M^{*}\cap(\ch(C)\setminus C)$; line~\ref{ln:fix} then computes
$\mathcal M^{*}\cap C$ by Lemma~\ref{lem:search}\ref{lemitem:fix}, and line~\ref{ln:decide}
accepts it and rejects the rest of $C$. Clause~\ref{D3} of $F$ is evaluated through
Lemma~\ref{lem:search}\ref{lemitem:de}.

For the bound let $m$ be that of Lemma~\ref{lem:finiteness}; as $M$
occupies at most $n$ input symbols, $\log m=O(n)+O(n^{2})\log M=\mathrm{poly}(n)$, and each of the at
most $2m$ rounds costs one decomposition of a subgraph of $\TabGraph$ together with at most $m$
applications of $F$, each a decomposition and $O(n)$ backward reachabilities linear in the subgraph,
so the run costs $\mathrm{poly}(m)=2^{\mathrm{poly}(n)}$.
\end{proof}




